% TOP_PROBLEM: {"id": 3087, "title": "Convex uniform 5-polytopes", "category_id": 6, "category": {"id": 6, "name": "geometry", "display_name": "Geometry", "description": "Euclidean and non-Euclidean geometry, geometric structures.", "slug": "geometry", "order_index": 6, "created_at": "2026-07-31T15:26:25.671Z"}, "status": "open", "problem_number": "OPG-37456", "set_id": 12}
% FIRST_POSTED: 2026-10-01
% ENTRY_KIND: statement_only
% UnsolvedMath ID 3087; source catalogue code OPG-37456
% !TeX program = lualatex
% Definition and sources only; this file is not an LLM solution attempt.
\documentclass[11pt,a4paper]{article}
\usepackage[margin=25mm]{geometry}
\usepackage{fontspec}
\setmainfont{lmroman10-regular.otf}[BoldFont=lmroman10-bold.otf,
 ItalicFont=lmroman10-italic.otf,BoldItalicFont=lmroman10-bolditalic.otf]
\usepackage{amsmath,amssymb,mathtools}
\usepackage{unicode-math}
\setmathfont{latinmodern-math.otf}
\AtBeginDocument{\let\setminus\smallsetminus}
\usepackage{xurl,hyperref}
\hypersetup{colorlinks=true,urlcolor=blue}
\setcounter{secnumdepth}{0}
\setlength{\parindent}{0pt}
\setlength{\parskip}{5pt}
\setlength{\emergencystretch}{3em}
\begin{document}
\section*{Convex uniform 5-polytopes}\label{3087}
{\small UnsolvedMath ID 3087 · OPG-37456 · Geometry · OpenGarden}
\subsection{Definitions and mathematical statement}
A convex five-dimensional polytope is the convex hull of finitely many points in $\mathbb R^5$, with affine span equal to $\mathbb R^5$. Its proper faces are intersections with supporting hyperplanes, and its facets are its four-dimensional faces. A Euclidean symmetry is an isometry of the ambient space carrying the polytope onto itself. Vertex transitivity means that, for every ordered pair of vertices, some such symmetry maps the first to the second.

Use the standard recursive notion of uniformity: a regular polygon is uniform in dimension two; in higher dimensions a convex polytope is uniform when its Euclidean symmetry group is vertex transitive and every facet is a uniform polytope. Segments provide the dimension-one base case. This allows different types of uniform facets in a single polytope; it does not require transitivity on all flags or on all faces of a fixed dimension.

The problem is to classify all convex uniform five-dimensional polytopes up to Euclidean similarity. A similarity is a composition of an isometry and multiplication of all distances by a common positive factor. Accordingly one may normalize every edge length to one. A complete enumeration must specify every infinite family and every exceptional type, identify repetitions under similarity, and prove that no additional convex uniform polytopes exist. Merely listing outputs of a particular construction is insufficient unless completeness for arbitrary uniform polytopes is proved.
\subsection{Short English statement}
Give and prove a complete classification of convex uniform five-dimensional polytopes, including all infinite families and exceptional shapes.
\subsection{Sources}
\begin{itemize}
\item[S1] ULAM.AI and the UnsolvedMath Contributors, supplied problems.json, record 3087 (OPG-37456), OpenGarden. Catalogue identity and source transcription; CC BY 4.0.
\url{https://huggingface.co/datasets/ulamai/UnsolvedMath}.
\item[S2] Open Problem Garden, Convex uniform 5-polytopes. Original problem and discussion.
\url{http://www.openproblemgarden.org/op/convex_uniform_5_polytopes}.
\item[S3] G. C. Shephard, A Construction for Wythoffian Polytopes, Canadian J. Math. 6 (1954), 133--144.
\url{https://doi.org/10.4153/CJM-1954-015-5}.
\end{itemize}
\end{document}
